\section*{Hoofdstuk \ref{ch:b2:frontier-orbitals} --- Grensorbitalen en reactiviteit}

\begin{solution}{exo:b2:frontier-orbitals:1}
Etheen: \omterm{def:b2:frontier-orbitals:frontier}{HOMO} $\alpha + \beta$, \omterm{def:b2:frontier-orbitals:frontier}{LUMO} $\alpha - \beta$. Allylanion (vier elektronen): \omterm{def:b2:frontier-orbitals:frontier}{HOMO} $\alpha$ (het
niet-bindende niveau), \omterm{def:b2:frontier-orbitals:frontier}{LUMO} $\alpha - 1.414\beta$. Butadieen: \omterm{def:b2:frontier-orbitals:frontier}{HOMO} $\alpha + 0.618\beta$, \omterm{def:b2:frontier-orbitals:frontier}{LUMO} $\alpha -
0.618\beta$.
\end{solution}

\begin{solution}{exo:b2:frontier-orbitals:2}
Hard: \ce{OH-}, \ce{H+}, \ce{F-}. Zacht: \ce{I-}, \ce{RS-}, het koolstofatoom van joodmethaan.
\end{solution}

\begin{solution}{exo:b2:frontier-orbitals:3}
Butadieen (s-cisconformeer bereikbaar) en cyclopentadieen (vastgezet in s-cis) reageren;
penta-1,4-dieen is niet geconjugeerd en reageert niet. Het (2\textit{Z},4\textit{Z})-dieen kan
de s-cisvorm alleen bereiken door zijn twee binnenste methylgroepen tegen elkaar te duwen: het
reageert zeer slecht.
\end{solution}

\begin{solution}{exo:b2:frontier-orbitals:4}
Cyclohex-3-een-1-carbonitril: een zesring, met de dubbele binding tussen de vroegere C2
en C3 van het \omterm{def:b2:frontier-orbitals:diels-alder}{dieen} en de groep \ce{C#N} op een van de twee ringkoolstofatomen die uit het
\omterm{def:b2:frontier-orbitals:diels-alder}{dienofiel} ontstaan zijn.
\end{solution}

\begin{solution}{exo:b2:frontier-orbitals:5}
$\Delta\varepsilon = \qty{8}{eV}$: storingsrekening $2\beta^2/\Delta\varepsilon = \qty{0.25}{eV}$. Exact: lager niveau
$-6 - \sqrt{16 + 1} = \qty{-10.123}{eV}$, winst $2 \times 0.123 = \qty{0.246}{eV}$.
\end{solution}

\begin{solution}{exo:b2:frontier-orbitals:6}
Joodmethaan is een \omterm{def:b2:frontier-orbitals:hard-soft}{zacht elektrofiel}: \omterm{def:b2:frontier-orbitals:control}{orbitaalcontrole}, aanval door het atoom met de grootste
HOMO-coëfficiënt, koolstof. Een vrij carbokation is een \omterm{def:b2:frontier-orbitals:hard-soft}{hard elektrofiel}: \omterm{def:b2:frontier-orbitals:control}{ladingscontrole},
gedeeltelijk aanval door stikstof, dat meer negatieve lading draagt.
\end{solution}

\begin{solution}{exo:b2:frontier-orbitals:7}
C4 van het \omterm{def:b2:frontier-orbitals:diels-alder}{dieen} (grootste HOMO-coëfficiënt) bindt aan het $\beta$-koolstofatoom van propenal (grootste
LUMO-coëfficiënt): 2-methoxycyclohex-3-een-1-carbaldehyde, met de twee substituenten op
naburige koolstofatomen (``ortho''-product).
\end{solution}

\begin{solution}{exo:b2:frontier-orbitals:8}
Een bicyclo[2.2.1]hepteenskelet (norborneen) met de estergroep op een koolstofatoom van het
vroegere \omterm{def:b2:frontier-orbitals:diels-alder}{dienofiel}, gericht naar de nieuwe \ce{C=C}-binding, onder de zesring
(endo), aan de kant tegenover de \ce{CH2}-brug.
\end{solution}

\begin{solution}{exo:b2:frontier-orbitals:9}
De twee carbonylgroepen verlagen de \omterm{def:b2:frontier-orbitals:frontier}{LUMO} van het \omterm{def:b2:frontier-orbitals:diels-alder}{dienofiel}: zijn afstand tot de \omterm{def:b2:frontier-orbitals:frontier}{HOMO} van
cyclopentadieen is veel kleiner dan die van etheen, de stabilisatie $\beta^2/\Delta\varepsilon$ van
de overgangstoestand is groter, en de barrière lager.
\end{solution}

\begin{solution}{exo:b2:frontier-orbitals:10}
De reactie is stereospecifiek: het maleaat geeft het adduct met beide estergroepen cis
(endo,endo als hoofdproduct; het is achiraal, een mesoverbinding); het fumaraat geeft het
trans-adduct, met één ester endo en één exo, gevormd als een paar enantiomeren (racemisch).
\end{solution}

\begin{solution}{exo:b2:frontier-orbitals:11}
Twee $\pi$-bindingen worden twee sterkere $\sigma$-bindingen: $\Delta_r H^\circ < 0$. Twee moleculen worden er één:
$\Delta_r S^\circ < 0$. $\Delta_r G^\circ = \Delta_r H^\circ - T\Delta_r S^\circ$ groeit met $T$ en wordt
positief boven $T_i = \Delta_r H^\circ/\Delta_r S^\circ$: bij hoge temperatuur splitst het adduct terug
(retro-diels-alder).
\end{solution}

\begin{solution}{exo:b2:frontier-orbitals:12}
Twee gevulde orbitalen van vrije elektronenparen naast elkaar vormen een wisselwerking met vier elektronen, die
afstoot. Het molecuul draait om de binding \ce{N-N} tot de vrije elektronenparen ongeveer
loodrecht op elkaar staan (een gaucheconformatie), waar hun overlap, en de afstoting,
klein is.
\end{solution}

\begin{solution}{pb:b2:frontier-orbitals:1}
\textbf{1.} Een vijfring met twee geconjugeerde dubbele bindingen en een \ce{CH2}: de ring
houdt de dieeneenheid in de s-cisconformatie.
\textbf{2.} Het \omterm{def:b2:frontier-orbitals:diels-alder}{dienofiel}, via een van zijn dubbele bindingen.
\textbf{3.} Een norborneenskelet dat gecondenseerd is met een cyclopenteenring; de twee nieuwe $\sigma$-bindingen verbinden
de uiteinden van het \omterm{def:b2:frontier-orbitals:diels-alder}{dieen} met de twee koolstofatomen van de dubbele binding van het \omterm{def:b2:frontier-orbitals:diels-alder}{dienofiel}.
\textbf{4.} Het \omterm{def:b2:frontier-orbitals:endo}{endo-adduct} (\omterm{def:b2:frontier-orbitals:endo}{endoregel}): de resterende ring van het \omterm{def:b2:frontier-orbitals:diels-alder}{dienofiel} ligt in de
overgangstoestand onder het \omterm{def:b2:frontier-orbitals:diels-alder}{dieen}.
\textbf{5.} Ja: ze is gecoördineerd, beide bindingen ontstaan aan dezelfde kant van elke partner, en de
geometrie van het \omterm{def:b2:frontier-orbitals:diels-alder}{dienofiel} blijft behouden.
\textbf{6.} Cyclopentadieen is zowel een goed \omterm{def:b2:frontier-orbitals:diels-alder}{dieen} (vastgezet in s-cis) als een \omterm{def:b2:frontier-orbitals:diels-alder}{dienofiel}, en de
reactie is gecoördineerd, zonder geladen intermediair dat gestabiliseerd moet worden.
\textbf{7.} Negatief: twee $\pi$-bindingen worden vervangen door twee $\sigma$-bindingen.
\textbf{8.} Negatief: twee moleculen geven er één.
\textbf{9.} $\Delta_r G^\circ = \Delta_r H^\circ - T\Delta_r S^\circ$ met beide termen negatief: negatief bij lage
$T$, nul bij $T_i = \Delta_r H^\circ/\Delta_r S^\circ$, positief daarboven.
\textbf{10.} Bij kamertemperatuur is $T < T_i$: dimerisatie is begunstigd, al is ze traag. Dicht bij
het kookpunt van het dimeer begunstigt het evenwicht het monomeer, en wordt het
snel bereikt.
\textbf{11.} Het monomeer verlaat het hete mengsel als damp: het verwijderen van een product laat het
evenwicht blijven verschuiven naar het monomeer (Le Chatelier).
\textbf{12.} Bij kamertemperatuur dimeriseert het opnieuw; koude vertraagt die reactie.
\textbf{13.} $0.62 - (-0.62) = 1.24|\beta|$.
\textbf{14.} $0.62 - (-0.35) = 0.97|\beta|$.
\textbf{15.} Cyclopentadieen met maleïnezuuranhydride: de kleinere afstand geeft een grotere
stabilisatie van de overgangstoestand, een snellere reactie.
\textbf{16.} Ze zijn geconjugeerd met de binding \ce{C=C} en elektronenzuigend: zoals in
propenal trekken ze het $\pi^*$-niveau omlaag.
\textbf{17.} Zijn \omterm{def:b2:frontier-orbitals:frontier}{LUMO}, die lobben heeft op de carbonylkoolstofatomen tegenover de centrale koolstofatomen van
de \omterm{def:b2:frontier-orbitals:frontier}{HOMO} van het \omterm{def:b2:frontier-orbitals:diels-alder}{dieen}.
\textbf{18.} Norborneenskelet; de anhydridering onder de nieuwe dubbele binding, tegenover de
\ce{CH2}-brug.
\textbf{19.} De twee bruggenhoofdkoolstofatomen en de twee koolstofatomen die het anhydride dragen: vier
stereocentra.
\textbf{20.} Nee: een spiegelvlak gaat door de \ce{CH2}-brug en tussen de twee helften
van het molecuul; het is een mesoverbinding.
\textbf{21.} Ze stonden cis op het \omterm{def:b2:frontier-orbitals:diels-alder}{dienofiel}, en de gecoördineerde additie houdt ze cis.
\textbf{22.} Hetzelfde skelet met de anhydridering aan de kant van de \ce{CH2}-brug.
\textbf{23.} Nee: van endo naar exo gaan zou betekenen dat \ce{C-C}-bindingen verbroken en opnieuw gevormd worden, via
de retroreactie, die verhitting vereist.
\textbf{24.} Water opent het anhydride: het cis-dicarbonzuur, met beide \ce{COOH}-groepen endo.
\textbf{25.} Het \omterm{def:b2:frontier-orbitals:endo}{endo-adduct} heeft $\boldsymbol{4}$ stereocentra.
\end{solution}
